\documentclass{beamer}

\usepackage[utf8]{inputenc}
\usepackage[spanish]{babel}
\usepackage{beamerthemeshadow}
\usepackage{times} %font times
\usepackage[T1]{fontenc} 
%\usepackage{apacite}
%\usepackage{fontspec}
%\usepackage[latin1]{inputenc}
\usepackage{multicol}
\usepackage[notocbib]{apacite}
\usepackage{tikz}
\usepackage{textpos}

%para que cuando se seleccione un texto las letras acentuadad y las � se copien bien
\usepackage{enumerate}
\usefonttheme{professionalfonts}

\addtobeamertemplate{frametitle}{}{%
	\begin{textblock*}{100mm}(0.9\textwidth,-1.08cm)
		\includegraphics[height=1cm,width=2cm]{figures/universidad-rey-juan-carlos}
\end{textblock*}}


%PARA LAS NOTAS
\usepackage{pgfpages}
%\setbeameroption{show notes on second screen}

\newcommand<>{\mynote}[1]{%
	\alt#2{%
		\note{\textcolor{red}{#1}}
	}
	{%
		\note{#1}
	}
}


\mode<presentation>{
\usetheme{Montpellier}
\setbeamercovered{transparent}
 \setbeamertemplate{navigation symbols}{}
 \usecolortheme{beaver}
 \setbeamercolor{local structure}{fg=darkred}
\setbeamercolor{structure}{fg=darkred}
}

\makeatletter
\defbeamertemplate*{footline}{Dan P theme}
{
  \leavevmode%
  \hbox{%
  \begin{beamercolorbox}[wd=.2\paperwidth,ht=2.25ex,dp=1ex,center]{author in head/foot}%
    \usebeamerfont{author in head/foot}URJC
  \end{beamercolorbox}%{}{~~(\insertshortinstitute)}
  \begin{beamercolorbox}[wd=.7\paperwidth,ht=2.25ex,dp=1ex,center]{title in head/foot}%
    \usebeamerfont{title in head/foot}\insertshorttitle
  \end{beamercolorbox}%
  \begin{beamercolorbox}[wd=.1\paperwidth,ht=2.25ex,dp=1ex,right]{date in head/foot}
\insertframenumber{} / \inserttotalframenumber\hspace*{2ex} 
  \end{beamercolorbox}}%
  \vskip0pt%
}
\makeatother

%Para la numeracion de los framebreak
\makeatletter
\defbeamertemplate*{frametitle continuation}{only if multiple}{%
	\ifnum \numexpr\beamer@endpageofframe+1-\beamer@startpageofframe\relax > 1
	\textmd{(%
		\insertcontinuationcountroman)}%
	\fi%
}
\makeatother

\setbeamertemplate{caption}{\raggedright\insertcaption\par} %eliminamos la palabra figura de las figuras

\newcommand\Wider[2][3em]{%
	\makebox[\linewidth][c]{%
		\begin{minipage}{\dimexpr\textwidth+#1\relax}
			\raggedright#2
		\end{minipage}%
	}%
}


\title{Uso de técnicas predictivas para la planificación de grupos en Secundaria y FP}
\vspace{-5mm}
\author{Universidad Rey Juan Carlos}

\institute{\textbf {Autor: Abel de Andrés Gómez\\ Tutor: Aurelio Berges García}} % auteur
\date{{\small 25 de junio de 2019}}




\begin{document} %inicio del documento

%portada
\begin{frame}[plain]{}
\begin{center}
\includegraphics [width =0.3 \textwidth ]{figures/escudo_urjc} %logo de la U en carpeta figures
\vspace*{-5mm}

\end{center}

\titlepage
\end{frame}

\mynote<2>{Buenas tardes a todos, soy Abel de Andres y con permiso de este tribunal paso a presentar mi TFM con el titulo de "Uso de tecnicas predictivas para la planificacion de grupos en secundaria y FP"}

%indice
\begin{frame}
\frametitle{Índice} %Esquema es el titulo de la diapositiva
	\tableofcontents
\end{frame}


\mynote<2>{%Este proyecto no se me ha ocurrido, en el año 2018 estuve trabajando con cargo a un convenio con la comunidad de Madrid. Contactaron conmigo a traves de una persona de la UPM (Marina). Estuve trabajando de junio a diciembre.En el convenio trabajaron 2 personas, 2 profesores y 3 personas de la Comunidad de Madrid. En ese convenio el objetivo era estudiar todos los datos de la comunidad de madrid para hacer una explotacion razonable. AL hilo de este proyecto, surge la necesidad de avanzar mas darle una vuelta de tuerca, surge la posibilidad de utilizar tecnicas predictivas para la gestion educativa. Estaba en el nucleo donde se toman las decisioines. Se trabajaba con el director del area tecnica. He interactuado con profesores que han sido directores y ahora son ATD (Asesor Tecnico Docente).
	En primer lugar y antes de nada, debo destacar que este proyecto no se me ha ocurrido a mi solo. En 2018 estuve trabajando en la Consejería de Educación a través de un convenio con la Universidad Politécnica de Madrid. Este convenio tenia por objeto estudiar y hacer una explotación razonables de los datos educativos de la Comunidad de Madrid. Para ello, he tenido que interactuar con profesores y directores que ahora son ATD(Asesor Tecnico docente).
	Una vez dicho esto paso a comentar los apartados en los que dividiremos esta presentación...
}

%\begin{frame}[plain]<beamer>{Outline}
%\tableofcontents[currentsection,currentsubsection]
%\end{frame}

%introducci�n

\section{Introducción} % Nueva Secci�n
\begin{frame}
\begin{center}
	\begin{beamercolorbox}[
		sep=8pt,center,rounded=true,shadow=true]{part title}
		\usebeamerfont{part title}
		\secname
	\end{beamercolorbox}
\end{center}
\end{frame}

\subsection*{Introducción} % Nueva Subsecci�n

\begin{frame}[allowframebreaks=1]
\frametitle{\secname : \subsecname}

\begin{block}{¿Que es la gestión educativa?} 
	Es un proceso continuo de toma de decisiones; desde las fases de planificación hasta las de evaluación, pasando por procesos
	intermedios de organización, asignación de recursos, implementación de procesos, etc. \cite{valenzuela2009construccion}
\end{block}

\begin{block}{¿Que es la Sobrepoblación en el aula?}  % define un marco
Es el exceso del número de estudiantes que se encuentran en un espacio determinado cuya capacidad no es adecuada para acogerlos ni cuenta con las condiciones adecuadas para el buen desenvolvimiento de los mismos. \cite{LILIA2013}
\end{block}

\begin{block}{¿Que es la ratio?}  % define un marco
Relación entre el número de alumnos y profesores, es un factor importante a la hora de realizar la planificación de los recursos y tiene gran influencia en la calidad educativa. \cite{TALIS2018}
\end{block} %acaba marco

\begin{itemize}
	\item La ratio creció en España en 2012 un 20\%, ahorrándose así 464 millones de euros.
	\item Las ratios en España son superiores a la media de la OCDE y la UE22. \cite{PANORAMA2018}
	\item Los recursos de las administraciones públicas no son infinitos.
	\item \textbf{¡Se debe planificar!}
	
	
\end{itemize}
	
\end{frame}



\subsection*{Objetivos}
\begin{frame}[allowframebreaks=1]
\frametitle{\secname : \subsecname}

El \textbf{objetivo general} es proponer un modelo para contribuir a la óptima planificación de los grupos escolares para los nuevos cursos, evitando así el gasto innecesario de recursos y controlando la sobrepoblación en el aula. \newline

\framebreak

Este objetivo se divide en los siguiente sub objetivos:

\begin{itemize}
	\item Seleccionar variables de interés, relativas a la resolución de la necesidad anteriormente expuesta por Unidad de Planificación, que aporten valor en el desarrollo de este TFM. 
	\item Estudiar la relación entre dichas variables con el propósito de comprender el contexto de la sobrepoblación en el aula y la planificación de grupos.
	\item Probar distintos modelos predictivos y seleccionar aquellos que aporten mayor precisión en la predicción de los grupos. 
	\item Obtener y utilizar el modelo de mayor precisión para realizar predicciones.
	
\end{itemize}
	
\end{frame}


\section{Justificación teórica}
\begin{frame}
\begin{center}
	\begin{beamercolorbox}[
		sep=8pt,center,rounded=true,shadow=true]{part title}
		\usebeamerfont{part title}
		\secname
	\end{beamercolorbox}
\end{center}
\end{frame}

\subsection{}
\begin{frame}
\frametitle{\secname}
\begin{block}{Objeto de búsqueda}
	Documentación científica acerca de la minería de datos en el ámbito educativo, concretamente, en la gestión de la educación.
	\begin{itemize}
		\item ScienceDirect, Scopus, Google Academics, etc.
	\end{itemize}
\end{block}

\begin{itemize}
	\item ¿Existen artículos relacionados con la minería de datos en la educación?
	\item ¿Existen artículos relacionados con la planificación educativa?
	\item ¿Que metodologías se siguen?
	\item ¿Que modelos predictivos contemplan?
	\item ¿Que variables se utilizan?
	\item ¿Que herramientas se usan?
\end{itemize}

\end{frame}

\section{Propuesta de intervención}
\begin{frame}
\begin{center}
	\begin{beamercolorbox}[
		sep=8pt,center,rounded=true,shadow=true]{part title}
		\usebeamerfont{part title}
		\secname
	\end{beamercolorbox}
\end{center}
\end{frame}

\subsection{}
\begin{frame}
\frametitle{\secname}

\makebox[\textwidth]{\includegraphics[width=12.5cm,height=6cm,keepaspectratio]{figures/PropuestaIntervencion3}}
\end{frame}

\section{Diseño de la investigación}

\begin{frame}
\begin{center}
	\begin{beamercolorbox}[
		sep=8pt,center,rounded=true,shadow=true]{part title}
		\usebeamerfont{part title}
		\secname
	\end{beamercolorbox}
\end{center}
\end{frame}

\begin{frame}
\frametitle{\secname : \subsecname}
\Wider[3em]{
	\begin{figure}[htb]
	\centering
	\caption{Fases del ciclo de vida de CRISP-DM. Recuperado de \protect\citeA{IBMCRISP2012}.}
	\includegraphics[width=0.5\textwidth,keepaspectratio]{../TemplateTFM/recursos/CRISPCicloIBM}
	\end{figure}
}




\end{frame}

\mynote<2>{
	\footnotesize
	\begin{itemize}
		\item Comprensión del negocio: Obtener información acerca de los objetivos de la investigación. Estudio de la situación de la educación en la Comunidad de Madrid.
			\item Comprensión de los datos:
			\item Preparación de los datos: Se realiza un tratamiento de los datos, se eliminan valores y variables vacías, se separan o se juntan variables a conveniencia de la investigación. Se separan los datos en dos subconjuntos (entrenamiento y pruebas)
			\item Modelado: Se utilizan y comparan los modelos predictivos. Para ello se utilizan métricas de selección de variables
			\item Evaluación: Se utilizan métricas de precisión para evaluar los modelos. Además, también se cuenta el tiempo que tardan en entrenarse. A partir de estas métricas se toma una decisión.
			\item Distribución: Toda la información recogida en la investigación se documenta. En esta parte se indica el modelo, las variables, herramientas y metodología que deben utilizarse para el control del proceso de la minería de datos.
	\end{itemize}
}

\section{Análisis de Resultados}
\begin{frame}
\begin{center}
	\begin{beamercolorbox}[
		sep=8pt,center,rounded=true,shadow=true]{part title}
		\usebeamerfont{part title}
		\secname
	\end{beamercolorbox}
\end{center}
\end{frame}


\subsection*{Análisis exploratorio}
\begin{frame}[allowframebreaks=1]
\footnotesize
\begin{block}{Análisis exploratorio}
	Se \textbf{examinan los datos} previo a la aplicación de técnicas estadísticas, para
	conseguir un entendimiento básico de estos y las relaciones existentes entre las variables analizadas. 
\end{block}
\frametitle{\secname : \subsecname}
 \begin{itemize}
 	\item Media de ratio en la CM -> 0,88; centros no sobrepoblados.
 	 \item Mayor sobrepoblación en DAT-Centro, seguido de DAT-Sur.
 	\item Número de alumnos y de grupos son importantes a la hora de predecir el número de grupos finales.
 	\item Variables mayor correladas con \textbf{ratio} son: número de alumnos, servicio de comedor, naturaleza del centro y código genérico del centro.
 	\item Variables mayor correladas con \textbf{número de grupos a planificar} son: número de alumnos, unidades, naturaleza, código genérico y servicio de comedor.
 \end{itemize}
 

\end{frame}
\subsection*{Análisis predictivo}
\begin{frame}
\begin{block}{Análisis predictivo}
	Se muestran los resultados obtenidos en el análisis de predicción.
\end{block}
\frametitle{\secname : \subsecname}
\begin{itemize}
	\item Las variables más importantes para realizar la predicción son: naturaleza del centro, número de curso, número de unidades, nivel de enseñanza y ratio.
	\item En la predicción: de 4436 datos existentes para grupos, únicamente 52 se modifican. De estos 52 grupos, 30 de ellos disminuyen en el número de unidades y 22 aumentan. Se observa para este curso 2017/2018, que \textbf{se reduce la población en el aula}.
\end{itemize}

\end{frame}

\section{Conclusiones}
\begin{frame}
\begin{center}
	\begin{beamercolorbox}[
		sep=8pt,center,rounded=true,shadow=true]{part title}
		\usebeamerfont{part title}
		\secname
	\end{beamercolorbox}
\end{center}
\end{frame}

\subsection*{Aportaciones del TFM}

\begin{frame}
\frametitle{\secname : \subsecname}
 \begin{itemize}
 	\item Uso de herramientas de minería de datos para la planificación académica.
 	\item Comparación de modelos predictivos en el entorno educativo para la planificación de grupos.
 	\item Descubrimiento de relaciones entre variables.
 	\begin{itemize}
 		\item Con resultados positivos para la gestión educativa.
 	\end{itemize}
 	\item Investigación relativa a la gestión y planificación en el entorno educativo.
 	\item Nuevas experiencias e ideas que se puedan iniciar no solo en la en la Consejería de Educación de la Comunidad de Madrid, sino en otras comunidades autónomas.
 \end{itemize}

\end{frame}

\subsection*{Conclusiones}

\begin{frame}[allowframebreaks=1]
\frametitle{\secname : \subsecname}
A continuación, se va a tratar de relacionar las conclusiones
con los objetivos establecidos:
\begin{enumerate}
	\item \textbf{Selección de variables:} Uso de algunas variables propuestas a partir de la instrucción de la Unidad de Planificación. \cite{INSTRCONSE}
	\begin{itemize}
		\item Se parte de 27 variables, finalmente se eligen 11. Se eliminan variables descriptivas (nombre del centro, código de este, etc.)
	\end{itemize}
	\item \textbf{Estudio de la relación entre variables para comprender el contexto de sobrepoblación en el aula} 
	\begin{itemize}
		\item Uso de la matriz de Correlación.
		\item Obtención de variables que hacen que la ratio aumente.
		\item Obtención de variables que mayor relación tengan con el número de grupos.
	\end{itemize}
	\item \textbf{Pruebas con distintos modelos predictivos y selección de los óptimos}
	\begin{itemize}
		\item Del total de modelos (6) se eligen 2: Árbol de decisión y Regresión Linear
		\item Mejores resultados (precisión y tiempo de entrenamiento).
	\end{itemize}
	\item \textbf{Se utiliza el algoritmo de mayor precisión en la predicción con datos existentes:} A partir del modelo de árbol de decisión se obtiene la predicción esperada para el curso 2017/2018 a partir del modelo entrenado a partir de los datos de 2016/2017.
\end{enumerate}
\framebreak
Por último:
\setbeamertemplate{items}[triangle]
\begin{itemize}
	\item Destacar la gran utilidad de este máster, que hace posible que los futuros docentes adquieran la capacidad de detectar y solucionar dificultades en el contexto educativo.
	\item Subrayar las competencias adquiridas con la realización de este TFM, relacionadas con el establecimiento de modelos de mejoras de la calidad en la educación y el analisis y conocimiento de la situación actual de esta.
\end{itemize}
%Destacar la gran utilidad de este máster en la formación de profesores, que hace posible que los futuros docentes adquieran la capacidad de detectar y solucionar cualquier dificultad existente en el contexto educativo para mejorar su calidad. No solamente dentro del aula, sino también en el plano de la gestión y planificación.


%\mynote<1>{note one\\}
%\mynote<2>{ note 2\\ }
%\mynote<3>{ note 3\\ }
%\mynote<1>{ note 4\\ }

\end{frame}



\subsection*{Líneas de Trabajo Futuro}

\begin{frame}
\frametitle{\secname : \subsecname}
\begin{itemize}
	\item Inclusión a la investigación de nuevas variables como aspectos físicos del centro, acceso a este, comunicaciones, estaciones de metro cercana, paradas de autobús, etc.
	\item Número y nuevas enseñanzas.
	\item Nuevos modelos de predicción y parámetros de ajustes.
	\item Realización de Software usando modelos investigados.
\end{itemize}
\end{frame}

\section{Referencias}
\begin{frame}
\begin{center}
	\begin{beamercolorbox}[
		sep=8pt,center,rounded=true,shadow=true,shadow=true]{part title}
		\usebeamerfont{part title}
		\secname
	\end{beamercolorbox}
\end{center}
\end{frame}

\begin{frame}[shrink=30]
\frametitle{Referencias}
	\bibliographystyle{apacite} % estilo de la bibliografia
	\bibliography{../TemplateTFM/referencias_prueba}

\end{frame}

\section*{FIN}
\begin{frame}{}
\centering \Huge
\emph{GRACIAS POR SU ATENCIÓN}
\end{frame}



\end{document}
